% !TEX program = xelatex
% 问题三分析——思维导图（紧凑竖排·箭距优化版，XeLaTeX 编译）
\documentclass[border=6pt]{standalone} % border=6pt：PDF 页面四周留白宽度
\usepackage{ctex}
\usepackage{fontspec}
\setmainfont{Times New Roman}
\usepackage{amsmath,amssymb}
\usepackage{tikz}
\usetikzlibrary{arrows.meta}

% ---------- 与前两问相同的配色 ----------
\definecolor{dgreen}{HTML}{0E4F43}  % 深绿边框/虚线
\definecolor{pOne}{HTML}{EAF6EC}    % Step1 面板底色（浅绿）
\definecolor{pTwo}{HTML}{E7F0FA}    % Step2 面板底色（浅蓝）
\definecolor{pThr}{HTML}{F6E9D4}    % Step3 面板底色（浅橙）
\definecolor{pFou}{HTML}{DEF2F7}    % Step4 面板底色（浅青）
\definecolor{pFiv}{HTML}{F8F1F7}    % Step5 面板底色（浅粉白）
\definecolor{pinkB}{HTML}{F7D8DC}   % 强调节点（粉红）
\definecolor{cyanA}{HTML}{8CDDE6}   % 青色节点
\definecolor{whiteA}{HTML}{F8F5FA}  % 公式框（浅紫白）
\definecolor{cyanC}{HTML}{DFF6F8}   % 结果节点（浅青）
\definecolor{greenC}{HTML}{EBF7ED}  % 方法节点（浅绿）

\begin{document}
	\begin{tikzpicture}[font=\normalsize, % 全图节点默认字号
		% --- 箭头样式：line width=1.6pt 线粗；length=6pt 头长；width=5pt 头宽 ---
		arr/.style={line width=1.6pt,-{Stealth[length=6pt,width=5pt]}},
		% --- 无箭头连接线（脊线用） ---
		wire/.style={line width=1.6pt},
		% --- 圆角矩形节点 ---
		boxR/.style={rounded corners=5pt,draw=dgreen,line width=2pt,
			inner sep=2.5mm,minimum height=10mm,align=center},
		% --- 直角矩形节点 ---
		boxS/.style={draw=dgreen,line width=2pt,
			inner sep=2.5mm,minimum height=10mm,align=center},
		% --- 主面板虚线边框 ---
		panel/.style={draw=dgreen,dotted,line width=2pt},
		]
		
		% ================= 五个虚线面板（左下、右上角坐标） =================
		\draw[panel,fill=pOne] (0,19.0) rectangle (24,22.0);  % Step1 面板：宽 24，高 3.0
		\draw[panel,fill=pTwo] (0,14.1) rectangle (24,18.6);  % Step2 面板：宽 24，高 4.5
		\draw[panel,fill=pThr] (0,7.7)  rectangle (24,13.7);  % Step3 面板：宽 24，高 6.0
		\draw[panel,fill=pFou] (0,2.8)  rectangle (24,7.3);   % Step4 面板：宽 24，高 4.5
		\draw[panel,fill=pFiv] (0,0)    rectangle (24,2.4);   % Step5 面板：宽 24，高 2.4
		
		% ================= Step 1：优化函数（箭距 1.0cm） =================
		\node[anchor=north west,font=\large] at (0.4,21.75)
		{\textbf{Step 1：优化函数——年均成本最小}};
		\node[boxS,fill=whiteA,font=\large] (J) at (7.6,20.1)
		{$\min\ J(X)=\dfrac{C_{\mathrm{purchase}}+c_{\mathrm{mid}}\,n_{\mathrm{mid}}+c_{\mathrm{big}}\,n_{\mathrm{big}}}{L}$};
		\node[boxR,fill=cyanC] (JP) at (17.2,20.1)
		{$C_{\mathrm{purchase}}=300$，$c_{\mathrm{mid}}=3$，$c_{\mathrm{big}}=12$（万元）\\$L$、$n_{\mathrm{mid}}$、$n_{\mathrm{big}}$ 由逐日仿真给出};
		\draw[arr] (J.east) -- (JP.west);
		
		% ================= Step 2：决策变量（决策向量框居中） =================
		\node[anchor=north west,font=\large] at (0.4,18.35) {\textbf{Step 2：决策变量}};
		\node[boxS,fill=cyanA] (X) at (12,17.1) {决策向量\quad $X=(T_{\mathrm{mid}},\ T_{\mathrm{big}},\ r_{\mathrm{th}})$};
		\node[boxS,fill=whiteA] (D1) at (4.2,15.2)  {$T_{\mathrm{mid}}$ 中维护间隔阈值\\固定间隔触发·$[40,120]$天};
		\node[boxS,fill=whiteA] (D2) at (12.0,15.2) {$T_{\mathrm{big}}$ 大维护间隔阈值\\固定间隔触发·$[150,400]$天};
		\node[boxS,fill=whiteA] (D3) at (19.8,15.2) {$r_{\mathrm{th}}$ 退化速率触发阈值\\状态响应触发·$[0.05,1.5]$};
		% 主干＋水平脊线 y=16.2＋三支竖直箭头
		\draw[wire] (X.south) -- (12,16.2);
		\draw[wire] (4.2,16.2) -- (19.8,16.2);
		\draw[arr] (4.2,16.2)  -- (D1.north);
		\draw[arr] (12.0,16.2) -- (D2.north);
		\draw[arr] (19.8,16.2) -- (D3.north);
		
		% ================= Step 3：约束条件（首行紧凑居中，间隙 0.8cm） =================
		\node[anchor=north west,font=\large] at (0.4,13.45) {\textbf{Step 3：约束条件}};
		\node[boxS,fill=greenC] (C1) at (7.3,12.0) {透水率为正\quad $P(t)>0,\ \forall t$};
		\node[boxS,fill=greenC] (C2) at (14.75,12.0)
		{年均下限\quad $\bar{P}_{365}(t)=\frac{1}{365}\displaystyle\sum_{i=0}^{364}P(t-i)\ge37,\ \forall t<L$};
		\node[boxS,fill=whiteA] (C3) at (12,10.55)
		{状态转移\quad $P(t+1)=P(t)-k_0\cdot\beta_{s(t)}\cdot(1+\alpha_{\mathrm{mid}})^{n_{\mathrm{mid}}(t)}\cdot(1+\alpha_{\mathrm{big}})^{n_{\mathrm{big}}(t)}+G(t)$};
		\node[boxS,fill=whiteA,font=\small] (C4) at (12,9.0)
		{维护触发\quad $\mathrm{do}_{\mathrm{big}}=(\Delta t_{\mathrm{big}}\ge T_{\mathrm{big}})$；\quad$\mathrm{do}_{\mathrm{mid}}=\neg\,\mathrm{do}_{\mathrm{big}}\wedge(\Delta t_{\mathrm{mid}}\ge T_{\mathrm{mid}}\ \vee\ \mathrm{rate}(t)\le -r_{\mathrm{th}})$};
		
		% ================= Step 4：遗传算法求解（箭距 0.6—0.8cm，新增搜索范围框） =================
		\node[anchor=north west,font=\large] at (0.4,7.05) {\textbf{Step 4：遗传算法求解}};
		\node[boxS,fill=pinkB,font=\small] (R) at (4.85,5.7) {约束非线性·维护致状态跳变\\目标函数不可解析 $\Rightarrow$ 选遗传算法};
		\node[boxS,fill=whiteA,font=\small] (GP) at (10.95,5.7)
		{GA参数：种群 30 · 代数 25\\交叉 0.7 · 变异 0.05};
		\node[boxS,fill=greenC,font=\small] (W) at (18.1,5.7)
		{搜索范围 $T_{\mathrm{mid}}\in[40,120]$·$T_{\mathrm{big}}\in[150,400]$\\$r_{\mathrm{th}}\in[0.05,1.5]$ 透水率/天};
		\draw[arr] (R.east) -- (GP.west);
		\draw[arr] (GP.east) -- (W.west);
		% 算法流程链（箭距 0.6cm）
		\node[boxS,fill=greenC] (G1) at (5.2,3.85)  {初始化种群};
		\node[boxS,fill=greenC] (G2) at (9.0,3.85)  {寿命仿真\\得 $L_i,\ n_{\mathrm{mid}},\ n_{\mathrm{big}}$};
		\node[boxS,fill=whiteA] (G3) at (12.4,3.85) {计算适应度 $J_i$};
		\node[boxS,fill=greenC] (G4) at (16.1,3.85) {选择·交叉·变异};
		\node[boxR,fill=cyanA] (G5) at (19.5,3.85) {输出 $X^{*}$};
		\draw[arr] (G1.east) -- (G2.west);
		\draw[arr] (G2.east) -- (G3.west);
		\draw[arr] (G3.east) -- (G4.west);
		\draw[arr] (G4.east) -- (G5.west);
		
		% ================= Step 5：结果对比（箭距 0.8cm） =================
		\node[anchor=north west,font=\large] at (0.4,2.15) {\textbf{Step 5：结果对比}};
		\node[boxS,fill=greenC] (S1) at (5.9,0.95) {优化后年均总成本 635.5 万元};
		\node[boxS,fill=cyanC] (S2) at (12.95,0.95) {平均年均成本 63.55 vs 原始 71.38（万元/年）};
		\node[boxR,fill=pinkB] (S3) at (19.05,0.9) {平均降幅 11.04\%\\9/10 台成本下降};
		\draw[arr] (S1.east) -- (S2.west);
		\draw[arr] (S2.east) -- (S3.west);
		
		% ================= 面板间流转箭头（端点恰在面板边线上，长 0.4cm） =================
		\draw[arr] (12,19.0) -- (12,18.6);  % Step1 → Step2
		\draw[arr] (12,14.1) -- (12,13.7);  % Step2 → Step3
		\draw[arr] (12,7.7)  -- (12,7.3);   % Step3 → Step4
		\draw[arr] (12,2.8)  -- (12,2.4);   % Step4 → Step5
		
	\end{tikzpicture}
\end{document}